\section{Data and methodology}
\label{sec:methodology}

\subsection{Microdata}

Our simulations use PolicyEngine UK \citep{policyengine}, an open-source static microsimulation model of the UK tax and benefit system, run on the Enhanced Family Resources Survey (EFRS) for 2023--24, uprated to fiscal year 2026--27. The EFRS augments the standard Family Resources Survey in two ways. First, because household surveys under-represent the top of the income distribution, synthetic high-income records are generated by training a quantile random forest on HMRC's Survey of Personal Incomes (SPI) and imputing the upper tail of incomes. Second, household weights are calibrated by gradient descent so that the weighted dataset matches a large set of national and regional administrative totals from the ONS, HMRC and DWP --- including population by age and region, income components, and benefit caseloads. Wealth variables, including property wealth and corporate (financial) wealth, are imputed from the ONS Wealth and Assets Survey (WAS).

\subsection{Council tax baseline}

The FRS records council tax imperfectly, so we reconstruct liabilities. Each household is assigned a council tax band (A--H in England and Scotland, A--I in Wales) and a bill equal to the average band rate in its region, with the 25 per cent single-person discount imputed to single-adult households. Band assignment probabilities are calibrated so that the distribution of dwellings by region and band matches Valuation Office Agency (and Scottish Assessors) counts. This bottom-up reconstruction yields, for 2026--27, gross council tax of \pounds 61.9 billion; netting off \pounds 4.3 billion of means-tested council tax reduction gives net revenue of
\begin{equation}
R_{\mathrm{CT}} \;=\; 61.9 - 4.3 \;=\; \pounds 57.6 \text{ billion},
\label{eq:ctrev}
\end{equation}
which is the amount the land value tax must replace. We note a calibration caveat: this model-based figure exceeds the Office for Budget Responsibility's accruals-based council tax receipts forecast of approximately \pounds 53.7 billion \citep{obr2025}. Were the replacement target instead calibrated to OBR-consistent receipts, the revenue-neutral rate would be approximately 0.71 rather than 0.77 per cent; all results below use the internally consistent model figure, which ensures that the simulated LVT replaces exactly the council tax the model abolishes household by household. Northern Ireland operates a separate domestic rates system, which we leave unchanged (see Section~\ref{sec:discussion}).

\subsection{Imputing land values}
\label{sec:landvalues}

The tax base is the unimproved value of land held directly (residential land under owner-occupied and other household-owned property) and indirectly (land on corporate balance sheets, attributed to households as ultimate owners). For household $i$ in region $r$, land value is
\begin{equation}
L_i \;=\; \underbrace{\lambda_r \, P_i}_{\text{household land}} \;+\; \underbrace{\frac{K_i}{\sum_j K_j w_j} \, L^{\mathrm{corp}}}_{\text{allocated corporate land}},
\label{eq:land}
\end{equation}
where $P_i$ is WAS-imputed property wealth, $\lambda_r$ is the region-specific land share of property value, $K_i$ is the household's corporate wealth, $w_j$ are survey weights, and $L^{\mathrm{corp}} = \pounds 2.06$ trillion is aggregate corporate-sector land. Land shares $\lambda_r$ vary from 85 per cent in London to 42 per cent in the North East, reflecting the fact that in high-price regions almost all of the price differential is location value rather than structure cost. Corporate land is allocated proportionally to each household's corporate wealth, treating shareholdings (direct and via pensions) as the ultimate claim on corporate assets.

Aggregates are calibrated to targets in the PolicyEngine UK data pipeline derived from the sectoral tables underlying the ONS National Balance Sheet \citep{ons2025}: UK land of \pounds 7.10 trillion at end-2024, split into \pounds 5.04 trillion household and \pounds 2.06 trillion corporate. (These derived targets differ from the headline bulletin presentation, which reports household non-produced assets of approximately \pounds 4.6 trillion in 2024 on a different sectoral grouping.) We uprate to 2026--27 using OBR per-capita nominal GDP growth \citep{obr2025}, giving a total land base of
\begin{equation}
L \;=\; \sum_i L_i \, w_i \;=\; \pounds 7.46 \text{ trillion},
\end{equation}
comprising \pounds 5.40 trillion of household land and \pounds 2.06 trillion of corporate land. The model total is 105.1 per cent of the (un-uprated) 2024 ONS benchmark. Land so measured is 68.4 per cent of total property wealth and 27.7 per cent of total household wealth (\pounds 26.98 trillion); mean land value per household is \pounds 232{,}672 and the median is \pounds 93{,}796.

\subsection{The reform and the revenue-neutral rate}

The reform (i) abolishes council tax in Great Britain and (ii) levies a flat annual tax at rate $r$ on each household's land value $L_i$, implemented in PolicyEngine UK via the parameters \texttt{gov.contrib.ubi\_center.land\_value\_tax.rate} and \texttt{gov.contrib.abolish\_council\_tax}. The budget-neutral rate equates LVT revenue with net council tax revenue:
\begin{equation}
r^{*} \;=\; \frac{R_{\mathrm{CT}}}{L} \;=\; \frac{57.6}{7{,}463} \;=\; 0.77\%.
\label{eq:rate}
\end{equation}
The net income change for household $i$ is $\Delta y_i = \mathrm{CT}_i - r^{*} L_i$, where $\mathrm{CT}_i$ is the abolished council tax liability (net of council tax reduction). Three maintained assumptions govern incidence, all standard in static microsimulation. First, incidence is assigned entirely to current landowners: we assume no behavioural response and no capitalisation of the tax into land prices. Second, rents are held fixed, so renters --- who own no residential land --- retain their abolished council tax liability in full and pay LVT only on any corporate wealth they hold. Third, the land base is held fixed across the rate grid, consistent with the inelastic supply of land. All results are for simulation year 2026, with poverty measured as absolute poverty against a baseline-fixed threshold, before housing costs (BHC) and after housing costs (AHC), and inequality by the Gini coefficient of equivalised household net income. In addition to the central rate $r^{*}=0.77\%$ we simulate a grid of rates $\{0.5, 1.0, 1.5, 2.0, 3.0, 5.0\}$ per cent to trace out revenue and distributional effects, and decompose revenue by base scope (household-only versus corporate-only land).
